\section{USAJMO 2026}
        \begin{enumerate}
        	\item (\textit{USAJMO 2026}) Let $a,b,c$ be distinct positive integers such that $ab+c=c^2$. Prove that $(a-b)^2\ge 4c$.


 

	\item (\textit{USAJMO 2026}) There are $m$ foxes and $n$ rabbits sitting in $m+n$ seats around a circular table, where all animals are distinguishable. If a fox and a rabbit are sitting next to each other, they may swap positions with each other. For a fixed starting configuration, determine the number of configurations that can be reached by a sequence of swaps.


 Note: Rotations and reflections of a configuration are considered distinct.


 

	\item (\textit{USAJMO 2026}) Let $ABC$ be an acute scalene triangle with no angle equal to $60^{\circ}$. Let $\omega$ be the circumcircle of $ABC$. Let $\triangle_B$ be the equilateral triangle with three vertices on $\omega$, one of which is $B$. Let $\ell_B$ be the line through the two vertices of $\triangle_B$ other than $B$. Let $\triangle_C$ and $\ell_C$ be defined analogously. Let $Y$ to be the intersection of $AC$ and $\ell_B$, and let $Z$ be the intersection of $AB$ and $\ell_C$.


 Suppose that the circumcircle of $AYZ$ intersects $\omega$ at $P\neq A$, $BC$ intersects $YZ$ at $D$, and $PA$ intersects $YZ$ at $E$. Prove that $PE=PD$.


 

	\item (\textit{USAJMO 2026}) Triangle $\triangle{ABC}$ has circumcircle $\omega$ and circumcenter $O$. Lines $AO$ and $BC$ meet at point $D$. Let $X$ be the $A$-excenter of $\triangle{ABD}$ and let $Y$ be the $A$-excenter of $\triangle{ACD}$. Prove that if $X$ lies on $\omega$, then $Y$ also lies on $\omega$.


 

	\item (\textit{USAJMO 2026}) A positive integer $n$ is called solitary if, for any nonnegative integers $a$ and $b$ such that $a+b=n$, either $a$ or $b$ contains the digit $1$ when written in base $10$. Determine, with proof, the number of solitary integers less than $10^{2026}$.


 

	\item (\textit{USAJMO 2026}) Emily has a red sheet of paper. She draws 2026 circles (not necessarily of equal size) on the piece of paper. She chooses a circle to color black, then cuts the paper around the circumference of all 2026 circles. She then separates the pieces of paper, into at least 2 black pieces and some number of red pieces. Is it possible that all black pieces are congruent?


 

\end{enumerate}